\documentclass[12pt,a4paper]{article}
\usepackage[utf8]{inputenc}
\usepackage[T1]{fontenc}
\usepackage{lmodern}
\usepackage[margin=2.5cm]{geometry}
\usepackage{setspace}
\onehalfspacing

\begin{document}

\noindent
Feira de Santana, April 2, 2026

\bigskip

\noindent
Dear Editor-in-Chief,\\
\textit{Stochastic Environmental Research and Risk Assessment}

\bigskip

\noindent
We are pleased to submit the manuscript entitled \textbf{``Analytical uncertainty of the statistical paradigm as a vector of stochastic risk in soil quality assessment under long-term conservation management''} for consideration as an original research article in \textit{Stochastic Environmental Research and Risk Assessment}.

\medskip

\noindent
Soil quality indices (SQI) are widely used to translate multivariate edaphic data into management recommendations, yet each study selects a single scoring algorithm and implicitly assumes that the resulting ranking is method-independent. Our study challenges this assumption by benchmarking five structurally distinct paradigms---ANOVA/GLM, fuzzy inference, PLS-SEM, MARS, and Random Forest---applied to the same long-term field experiment (22 years, $n = 108$, 18 indicators across four edaphic domains on a Typic Hapludult in northeast Brazil). We quantified inter-paradigm agreement using formal concordance metrics (Kendall~$W$, ICC, Bland--Altman limits of agreement), Borda count aggregation, and Monte Carlo simulations of rank stability.

\medskip

\noindent
The study has direct relevance to the scope of SERR\null A in three respects:

\begin{enumerate}
  \item \textbf{Stochastic risk quantification.} Monte Carlo resampling demonstrated that concordance between linear and non-linear families stabilizes only above 30 subsamples, and Bland--Altman analysis revealed that linear-only approaches may yield practical disagreement for intermediate management positions, introducing stochastic risk into conservation recommendations.
  \item \textbf{Multi-paradigm uncertainty assessment.} By applying concordance theory across paradigms of unequal structural complexity, the study frames soil quality assessment as a metrological problem of analytical uncertainty propagation rather than a simple scoring exercise.
  \item \textbf{Decision-support under model-choice uncertainty.} The results indicate that non-linear methods (fuzzy, MARS, Random Forest) form a block of near-interchangeable rankings, whereas the exclusive use of linear models may compromise recommendation reliability, a finding with direct implications for environmental risk management.
\end{enumerate}

\medskip

\noindent
To the best of our knowledge, no previous study has simultaneously compared five paradigms from different analytical families on the same edaphic dataset using formal stochastic concordance metrics. The manuscript therefore fills a relevant gap at the interface of pedometrics, environmental monitoring, and risk assessment.

\medskip

\noindent
The manuscript has not been published or submitted elsewhere. All authors have read and approved the final version. We declare no conflicts of interest.

\bigskip

\noindent
We look forward to your consideration.

\bigskip

\noindent
Sincerely,

\bigskip\bigskip

\noindent
Luiz Diego Vidal Santos\\
(Corresponding author)\\
DCHF, Universidade Estadual de Feira de Santana\\
Feira de Santana, Bahia, Brasil\\
E-mail: ldvsantos@uefs.br

\end{document}
